\subsection{Primary submodules corresponding to graded prime ideals}

PROPOSITION 3. Let \(\Delta\) be a torsion-free commutative group, A \(a\) graded Noetherian ring of type \(\Delta\) , \(\mathfrak{p}\) a graded ideal of A and M \(a\) graded A-module of type \(\Delta\) not reduced to 0. Suppose that for every homogeneous element \(a\) of \(\mathfrak{p}\) the homothety of ratio \(a\) on M is almost nilpotent and that for every homogeneous element \(b\) of A - \(\mathfrak{p}\) the homothety of ratio \(b\) on M is injective. Then \(\mathfrak{p}\) is prime and the submodule \(\{0\}\) of M is \(\mathfrak{p}\) -primary.

It suffices to show that \(\operatorname{Ass}(\mathbf{M}) = \{\mathfrak{p}\}\) ( \(\S 2\) , no. 1, Proposition 1). Let \(\mathfrak{q}\) be a prime ideal associated with \(\mathbf{M}\) ; it is a graded ideal and it is the annihilator of a homogeneous element \(x \neq 0\) of \(\mathbf{M}\) (no. 1, Proposition 1). For every homogeneous element \(a\) of \(\mathfrak{q}\) , \(ax = 0\) and hence the homothety of ratio \(a\) on \(\mathbf{M}\) is not injective, whence \(a \in \mathfrak{p}\) . Conversely, let \(b\) be a homogeneous element of \(\mathfrak{p}\) ; there exists an integer \(n > 0\) such that \(b^n x = 0\) , whence \(b^n \in \operatorname{Ann}(x) = \mathfrak{q}\) and, as \(\mathfrak{q}\) is prime, \(b \in \mathfrak{q}\) . As \(\mathfrak{p}\) and \(\mathfrak{q}\) are generated by their respective homogeneous element, \(\mathfrak{p} = \mathfrak{q}\) , which proves that \(\operatorname{Ass}(\mathbf{M}) \subset \{\mathfrak{p}\}\) . As \(\mathbf{M} \neq \{0\}\) , \(\operatorname{Ass}(\mathbf{M}) \neq \emptyset\) ( \(\S 1\) , no. 1, Corollary 1 to Proposition 2), whence \(\operatorname{Ass}(\mathbf{M}) = \{\mathfrak{p}\}\) .

PROPOSITION 4. Let \(\Delta\) be a torsion-free commutative group, A a graded Noetherian ring of type \(\Delta\) and M a graded A-module of type \(\Delta\) . Let \(\mathfrak{p}\) be a prime ideal of A and N, a submodule of M which is \(\mathfrak{p}\) -primary with respect to M.

\begin{enumerate}
\def\labelenumi{(\roman{enumi})}
\item
  The largest graded ideal \(\mathfrak{p}'\) of A contained in \(\mathfrak{p}\) (Algebra, Chapter II, § 11, no. 3) is prime.
\item
  The largest graded submodule \(\mathbf{N}'\) of \(\mathbf{N}\) is \(\mathfrak{p}'\) -primary with respect to \(\mathbf{M}\) .
\end{enumerate}

We know (loc. cit.) that the homogeneous elements of \(p'\) (resp. \(N'\) ) are just the homogeneous elements of p (resp. N). Let a be a homogeneous element of p; if x is a homogeneous element of M, there exists an integer n \textgreater{} 0 such that \(a^{n}x \in N\) ; as \(a^{n}x\) is homogeneous, \(a^{n}x \in N'\) ; as every \(y \in M\) is the direct sum of a finite number of homogeneous elements, we conclude that there exists an integer q \textgreater{} 0 such that \(a^{q}y \in N'\) , so that the homothety with ratio a in \(M/N'\) is almost nilpotent.

Consider now a homogeneous element \(b\) of \(A - p'\) ; then \(b \notin p\) since \(b\) is homogeneous. Let \(x\) be an element of \(M\) such that \(bx \in N'\) and let \((x_i)_{i \in \Delta}\) be the family of homogeneous components of \(x\) . As \(N'\) is graded, \(bx_i \in N'\) for all \(i\) , hence \(bx_i \in N\) and, as \(b \notin p\) , we conclude that \(x_i \in N\) ; as \(x_i\) is homogeneous, \(x_i \in N'\) , whence \(x \in N'\) and the homothety with ratio \(b\) on \(M/N'\) is injective. Proposition 4 then follows from Proposition 3 applied to \(p'\) and \(M/N'\) .

